% Generated draft: requires mathematical and source review.
\documentclass[11pt]{article}
\usepackage[margin=25mm]{geometry}
\usepackage{fontspec}
\setmainfont{Latin Modern Roman}
\usepackage{amsmath,amssymb,mathtools,unicode-math}
\setmathfont{Latin Modern Math}
\usepackage{enumitem,needspace,xurl,hyperref}
\setlength{\parindent}{0pt}
\setlength{\parskip}{5pt}
\setlength{\emergencystretch}{3em}
\newcommand{\N}{\mathbb{N}}
\newcommand{\Z}{\mathbb{Z}}
\newcommand{\Q}{\mathbb{Q}}
\newcommand{\R}{\mathbb{R}}
\newcommand{\C}{\mathbb{C}}
\newcommand{\F}{\mathbb{F}}
\newcommand{\A}{\mathbb{A}}
\newcommand{\PP}{\mathbb{P}}
\newcommand{\E}{\mathbb{E}}
\newcommand{\eps}{\varepsilon}
\newcommand{\dd}{\,\mathrm{d}}
\begin{document}
\section*{UnsolvedMath ID: 30001732}
\textbf{Source problem code:} catalogue entry
\subsection{Definitions and mathematical statement}
The supplied record uses \(D_\kappa(x,F)\) to denote kernelized spatial depth, with \(x\) the evaluation argument, \(F\) the distribution argument, and \(\kappa\) the kernel. It supplies no defining formula, domain for \(x\), admissible class of distributions \(F\), admissible class of kernels, or parameterization of the kernel. These roles describe the notation only; they do not provide a complete mathematical definition.

The source poses the following research questions:
\begin{enumerate}
\item What properties does \(D_\kappa(x,F)\) possess?
\item How should the kernel parameter be chosen optimally?
\item What is the relationship between \(D_\kappa(x,F)\) and kernel density estimation?
\end{enumerate}
The formulation remains unresolved at the level of hypotheses and targets. The properties to be established are not specified. The parameter space, optimality criterion, statistical setting, and dependence of the choice on \(F\), \(x\), or observed data are not stated. The kernel density estimator and the intended meaning of its relationship to depth are also unspecified. Consequently, no precise universally quantified assertion, optimization problem, or comparison theorem can be extracted from this record alone. Definitions, admissibility and existence conditions, and any relevant boundary or limiting regimes must be recovered from the cited report before a more specific mathematical statement is justified.
\subsection{Short English statement}
The source asks which properties kernelized spatial depth has, how to choose its kernel parameter optimally, and how it relates to kernel density estimation. The supplied record lacks the definitions, hypotheses, and comparison or optimization criteria needed to turn these questions into a precise mathematical conjecture.
\subsection{Sources}
\begin{itemize}
\item Supplied local UnsolvedMath registry record, permanent ID \(30001732\), source problem code \texttt{catalogue entry}, titled \emph{Kernel Selection and Properties of Kernelized Spatial Depth}, in the Oberwolfach Reports Open Problems collection. The record gives the proposed year as 2011 and identifies no individual proposer.
\item Supplied report citation: \emph{Mini-Workshop: Level Sets and Depth Contours in High Dimensional Data} (2011), Oberwolfach Reports, pp.\ 691--717, DOI \texttt{10.4171/owr/2011/13}. Source URL and supplied literature-source link, labeled ``Oberwolfach Report 13/2011'': \url{https://doi.org/10.4171/owr/2011/13}. The supplied background describes Oberwolfach Reports as records of mathematical workshops and their open-problem discussions.
\item Supplied literature assessment, checked on exactly \texttt{2026-08-21}: the entry is labeled open, and the assessment reports that no complete general answers to the extracted questions were located. This is supplied dated curation derived from the OWR material, not an independently verified literature assessment; its status and cited claims require independent verification.
\item The registry labels the statement \texttt{corrected\_verified} and reports that its mathematical content was checked against the official report, with notation, context, and question form normalized without changing the source-level claim. This is supplied verification metadata, not verification performed for this draft.
\end{itemize}
\end{document}
